\documentclass[12pt,a4paper]{article}

\usepackage[utf8]{inputenc}
\usepackage[T1]{fontenc}
\usepackage[brazil]{babel}
\usepackage{graphicx}
\usepackage{hyperref}
\usepackage{listings}
\usepackage{xcolor}
\usepackage{geometry}
\usepackage{indentfirst}

\geometry{margin=3cm}

\definecolor{codebg}{RGB}{245,245,245}
\lstset{
  backgroundcolor=\color{codebg},
  basicstyle=\ttfamily\small,
  breaklines=true,
  frame=single,
  numbers=left,
  numberstyle=\tiny,
  showstringspaces=false,
  tabsize=2
}

\title{Relat\'orio das Atividades Pr\'aticas de Fundamentos de DevOps}
\author{NOME DO ALUNO}
\date{\today}

\begin{document}

\maketitle

\tableofcontents
\newpage

\section{Introdu\c{c}\~ao}

DevOps \'e uma cultura e um conjunto de pr\'aticas que aproxima as equipes de
desenvolvimento (Dev) e de opera\c{c}\~oes (Ops), com o objetivo de entregar
software com mais qualidade, frequ\^encia e seguran\c{c}a. A integra\c{c}\~ao entre
desenvolvimento e opera\c{c}\~oes \'e importante porque reduz o tempo entre a
cria\c{c}\~ao de uma funcionalidade e sua disponibiliza\c{c}\~ao ao usu\'ario,
diminuindo falhas de comunica\c{c}\~ao e aumentando a confiabilidade do ambiente.

A automa\c{c}\~ao, o versionamento, a Infraestrutura como C\'odigo (IaC), os
containers, a an\'alise de c\'odigo e os pipelines s\~ao importantes porque tornam
o processo repet\'ivel, rastre\'avel e menos sujeito a erros manuais. Com essas
pr\'aticas, o ambiente pode ser recriado a qualquer momento, as mudan\c{c}as ficam
registradas e a qualidade do c\'odigo pode ser verificada antes do deploy.

O objetivo das pr\'aticas realizadas na disciplina foi construir um ambiente
completo de laborat\'orio, passando pela cria\c{c}\~ao de m\'aquinas virtuais com
Vagrant, pelo versionamento com Git e GitHub, pela automa\c{c}\~ao com Ansible e
pela integra\c{c}\~ao com SonarQube e Jenkins dentro do ciclo DevOps.

\section{Desenvolvimento}

\subsection{Atividade com Vagrant}

O Vagrant \'e uma ferramenta de automa\c{c}\~ao de ambientes que permite criar e
gerenciar m\'aquinas virtuais de forma reproduz\'ivel a partir de um arquivo de
configura\c{c}\~ao chamado \texttt{Vagrantfile}. Na disciplina, ele foi utilizado
para criar um laborat\'orio com tr\^es m\'aquinas virtuais: um n\'o de controle
(\texttt{control-node}), um servidor de aplica\c{c}\~ao (\texttt{app01}) e um
servidor de banco de dados (\texttt{db01}).

\begin{itemize}
  \item \textbf{Box utilizada:} \texttt{generic/rocky9}
  \item \textbf{Provider:} libvirt (KVM)
  \item \textbf{M\'aquinas:} 3 VMs
  \item \textbf{Hostnames e IPs:} \texttt{control-node} (192.168.1.2),
        \texttt{app01} (192.168.1.3) e \texttt{db01} (192.168.1.4)
  \item \textbf{Portas redirecionadas:} 80 do \texttt{app01} para 8080 do host;
        3306 do \texttt{db01} para 3306 do host
\end{itemize}

O \texttt{Vagrantfile} principal define a box, o nome de cada m\'aquina, o IP da
rede privada, as portas redirecionadas, o provisionamento via shell script e a
quantidade de mem\'oria e CPUs de cada VM.

\begin{lstlisting}[language=Ruby,caption={Trecho do Vagrantfile}]
Vagrant.configure("2") do |config|
  config.vm.box = "generic/rocky9"

  config.vm.define "control-node" do |c|
    c.vm.hostname = "control-node"
    c.vm.network "private_network", ip: "192.168.1.2"
    c.vm.provision "shell", path: "control-node/provision.sh"
    c.vm.provider "libvirt" do |libvirt|
      libvirt.memory = 2048
      libvirt.cpus = 2
    end
  end
end
\end{lstlisting}

Comandos utilizados:

\begin{lstlisting}[language=bash,caption={Comandos do Vagrant}]
vagrant up
vagrant status
vagrant ssh control-node
vagrant halt
vagrant reload
vagrant destroy
\end{lstlisting}

\begin{figure}[htbp]
  \centering
  % TODO: substituir pelo print do vagrant status com as maquinas em execucao
  % \includegraphics[width=0.9\textwidth]{imagens/vagrant-status.png}
  \caption{Execu\c{c}\~ao do \texttt{vagrant status} com as m\'aquinas ativas.}
\end{figure}

\subsection{Atividade com Git e GitHub}

O Git \'e um sistema de controle de vers\~ao distribu\'ido que registra o
hist\'orico de altera\c{c}\~oes de um projeto. Um reposit\'orio local fica na
m\'aquina do desenvolvedor, enquanto o reposit\'orio remoto fica hospedado em um
servi\c{c}o como o GitHub. O \textit{commit} registra um conjunto de altera\c{c}\~oes
no hist\'orico; o \textit{push} envia os commits para o remoto; o \textit{pull}
traz as altera\c{c}\~oes do remoto; e a \textit{branch} permite trabalhar em
ramifica\c{c}\~oes paralelas do projeto.

Na disciplina, o Git foi utilizado para versionar os arquivos do ambiente
(Vagrantfile, scripts de provisionamento, invent\'ario e playbooks do Ansible,
Jenkinsfile e o c\'odigo da aplica\c{c}\~ao). O versionamento \'e importante em
DevOps porque permite rastrear mudan\c{c}as, voltar a vers\~oes anteriores e
integrar o c\'odigo a pipelines de CI/CD.

\begin{lstlisting}[language=bash,caption={Comandos do Git utilizados}]
git init
git status
git add .
git commit -m "infraestrutura devops"
git remote add origin https://github.com/USUARIO/trabalho-devops.git
git push -u origin main
git branch
git checkout -b nova-branch
git merge nova-branch
\end{lstlisting}

\begin{figure}[htbp]
  \centering
  % TODO: print do repositorio no GitHub
  % \includegraphics[width=0.9\textwidth]{imagens/github-repositorio.png}
  \caption{Reposit\'orio do projeto no GitHub.}
\end{figure}

\begin{figure}[htbp]
  \centering
  % TODO: print do git status e do git log
  % \includegraphics[width=0.9\textwidth]{imagens/git-status.png}
  \caption{Comandos \texttt{git status} e \texttt{git log} executados.}
\end{figure}

\subsection{Atividade com Ansible}

O Ansible \'e uma ferramenta de automa\c{c}\~ao de configura\c{c}\~ao que permite
descrever o estado desejado dos servidores em arquivos YAML chamados playbooks.
Ele implementa o conceito de Infraestrutura como C\'odigo, pois a configura\c{c}\~ao
do ambiente fica registrada em arquivos version\'aveis. O \textit{Control Node} \'e
a m\'aquina que executa o Ansible; os \textit{Managed Hosts} s\~ao os servidores
gerenciados. A comunica\c{c}\~ao \'e feita via SSH, sem necessidade de agente
instalado nos hosts.

O invent\'ario foi configurado em \texttt{/home/vagrant/ansible/inventory.ini},
com o grupo \texttt{[control]} e o grupo \texttt{[servidores]}:

\begin{lstlisting}[caption={Invent\'ario do Ansible}]
[control]
control-node ansible_host=192.168.1.2

[servidores]
app01 ansible_host=192.168.1.3
db01 ansible_host=192.168.1.4

[servidores:vars]
ansible_user=vagrant
\end{lstlisting}

O playbook \texttt{app.yml} instala e inicia o \texttt{httpd} no \texttt{app01},
cria a p\'agina inicial e libera a porta 80 no firewall. O playbook \texttt{db.yml}
instala e inicia o \texttt{mariadb-server} no \texttt{db01} e libera a porta 3306.

\begin{lstlisting}[language=YAML,caption={Trecho do playbook app.yml}]
---
- name: Configurar servidor de aplicacao
  hosts: app01
  become: yes

  tasks:
    - name: Instalar o httpd
      dnf:
        name: httpd
        state: present

    - name: Habilitar e iniciar o httpd
      service:
        name: httpd
        state: started
        enabled: yes
\end{lstlisting}

\begin{lstlisting}[language=bash,caption={Comandos do Ansible}]
ansible all -m ping
ansible-playbook playbooks/site.yml
ssh vagrant@192.168.1.3
\end{lstlisting}

\begin{figure}[htbp]
  \centering
  % TODO: print do ansible all -m ping
  % \includegraphics[width=0.9\textwidth]{imagens/ansible-ping.png}
  \caption{Teste de conectividade com \texttt{ansible all -m ping}.}
\end{figure}

\begin{figure}[htbp]
  \centering
  % TODO: print da execucao do playbook
  % \includegraphics[width=0.9\textwidth]{imagens/ansible-playbook.png}
  \caption{Execu\c{c}\~ao do playbook \texttt{site.yml}.}
\end{figure}

\subsection{Atividade com SonarQube}

O SonarQube \'e uma plataforma de an\'alise est\'atica de c\'odigo que identifica
bugs, vulnerabilidades, code smells e duplica\c{c}\~oes. A an\'alise est\'atica
\'e importante porque encontra problemas antes que o c\'odigo chegue ao ambiente
de produ\c{c}\~ao, contribuindo para a qualidade e a seguran\c{c}a da aplica\c{c}\~ao.

O SonarQube foi executado em uma m\'aquina virtual dedicada
(\texttt{sonarqube}, IP 192.168.1.5), via container Podman, com acesso pelo
navegador em \texttt{http://localhost:9000}. O projeto foi criado no painel e a
an\'alise foi executada com o \texttt{sonar-scanner}.

\begin{lstlisting}[language=bash,caption={Comando do sonar-scanner}]
sonar-scanner \
  -Dsonar.projectKey=app-exemplo \
  -Dsonar.sources=. \
  -Dsonar.host.url=http://192.168.1.5:9000 \
  -Dsonar.login=TOKEN_OCULTADO
\end{lstlisting}

\begin{figure}[htbp]
  \centering
  % TODO: print do dashboard do SonarQube com os resultados
  % \includegraphics[width=0.9\textwidth]{imagens/sonarqube-dashboard.png}
  \caption{Dashboard do SonarQube com o resultado da an\'alise.}
\end{figure}

\subsection{Atividade com Jenkins}

O Jenkins \'e um servidor de automa\c{c}\~ao usado para implementar pr\'aticas de
CI/CD. A Integra\c{c}\~ao Cont\'inua (CI) automatiza o build e os testes a cada
altera\c{c}\~ao no c\'odigo; a Entrega Cont\'inua garante que o software esteja
sempre pronto para ser implantado; e o Deploy Cont\'inuo automatiza a implanta\c{c}\~ao
em produ\c{c}\~ao.

O Jenkins foi executado em uma m\'aquina virtual dedicada (\texttt{jenkins}, IP
192.168.1.6), via container Podman, com acesso em \texttt{http://localhost:8081}.
Foi criado um pipeline declarativo com as etapas de Checkout, Build, Test,
An\'alise SonarQube e Deploy.

\begin{lstlisting}[language=Groovy,caption={Trecho do Jenkinsfile}]
pipeline {
    agent any

    stages {
        stage('Build') {
            steps {
                echo 'Executando build da aplicacao'
            }
        }
        stage('Test') {
            steps {
                echo 'Executando testes'
            }
        }
        stage('Deploy') {
            steps {
                echo 'Executando deploy'
            }
        }
    }
}
\end{lstlisting}

\begin{figure}[htbp]
  \centering
  % TODO: print do pipeline executado com sucesso
  % \includegraphics[width=0.9\textwidth]{imagens/jenkins-pipeline.png}
  \caption{Pipeline executado com sucesso no Jenkins.}
\end{figure}

\section{Conclus\~ao}

As pr\'aticas permitiram compreender como as ferramentas se integram dentro do
ciclo DevOps: o Vagrant cria o ambiente, o Git versiona a infraestrutura e o
c\'odigo, o Ansible automatiza a configura\c{c}\~ao, o SonarQube verifica a
qualidade e o Jenkins automatiza o pipeline. O resultado \'e um processo mais
r\'apido, confi\'avel e reproduz\'ivel.

\section{Pergunta Reflexiva Final}

% TODO: escrever a resposta com as suas proprias palavras (minimo 10 linhas),
% abordando: tecnologia nova/pouco conhecida, o que aprendeu, maior dificuldade,
% como a pratica ajudou, como as ferramentas se conectam e a importancia para
% a formacao profissional.

\section*{Refer\^encias}

\begin{itemize}
  \item Documenta\c{c}\~ao oficial do Vagrant: \url{https://developer.hashicorp.com/vagrant/docs}
  \item Documenta\c{c}\~ao oficial do Git: \url{https://git-scm.com/doc}
  \item Documenta\c{c}\~ao oficial do Ansible: \url{https://docs.ansible.com/}
  \item Documenta\c{c}\~ao oficial do SonarQube: \url{https://docs.sonarsource.com/sonarqube/}
  \item Documenta\c{c}\~ao oficial do Jenkins: \url{https://www.jenkins.io/doc/}
\end{itemize}

\end{document}
